\section{Engine Architecture}
\subsection{Design goals}
Four constraints derived from educational deployment govern the design. Scenario authors must interact exclusively with high-level abstractions and never encounter Three.js types, lowering the graphics-programming threshold for content teams. The API must follow Unity conventions in naming, lifecycle and semantics, leveraging the pool of developers already familiar with that engine. The engine must remain lightweight, since bundle size and initialization time are first-class concerns for classroom hardware on institutional networks. And runtime performance, particularly on the integrated graphics typical of student laptops, must suffice for stable interactive use.

\subsection{High-level architecture}

The engine is organized into five layers, shown in Fig.~\ref{fig:architecture}: user-authored scenarios, the public engine API, runtime subsystems, an internal Three.js adapter, and the hidden rendering backend. This separation follows the subsystem stratification found in mature game-engine architectures~\cite{Ullmann2025}.

\texttt{Application} manages the engine lifetime, including WebGL context creation, the \texttt{requestAnimationFrame} loop and subsystem initialization. \texttt{SceneManager} maintains the active scene graph and lifecycle events, while \texttt{Resources}, input and physics provide runtime services. The public API exposes Unity-inspired types such as \texttt{GameObject}, \texttt{Component}, \texttt{Transform} and \texttt{Material}. Engine state is synchronized with Three.js only through the internal adapter layer, keeping the rendering backend hidden from scenario code.

\begin{figure}[!htb]
\centering
\fbox{\parbox{0.9\linewidth}{\centering
\includegraphics[width=0.8\linewidth]{fig1_architecture.png}}}
\caption{Engine layered architecture.}
\label{fig:architecture}
\end{figure}
\FloatBarrier

\subsection{Core subsystems}
Logic is composed from \texttt{Component} instances attached to \texttt{GameObject} containers, mirroring Unity; \texttt{ScriptableBehaviour} adds serialization, so components can be authored, instantiated from data and persisted across scene reloads. Composition rests on a hierarchical \texttt{Transform} graph exposing the standard Unity properties (\texttt{localPosition}, \texttt{localRotation}, \texttt{localScale}, \texttt{position}, \texttt{rotation}, \texttt{forward}) and methods (\texttt{translate}, \texttt{rotate}, \texttt{lookAt}, \texttt{transformPoint}). That hierarchy parallels the underlying Three.js \texttt{Object3D} graph but is conceptually independent -- a separation aligned with engine-level geometric optimization patterns reported for AR pipelines, where pre-processing of scene-graph elements is essential to interactive performance on consumer devices~\cite{Aung2022}.

Materials follow Unity's \texttt{sharedMaterial} versus \texttt{material} distinction: assigning \texttt{sharedMaterial} produces a reference shared across instances, while accessing \texttt{material} triggers an automatic clone for per-instance editing. \texttt{Standard\allowbreak Material} (PBR), \texttt{UnlitMaterial} and a low-level \texttt{Material} for custom shaders are provided. Assets load asynchronously from the active archive through \texttt{Resources.\allowbreak tryLoad<T>(type, path)}, with decoders for textures, JSON, text, audio and 3D models. Content is packaged as a ZIP archive of manifest, compiled scripts and asset folders; once loading completes the scenario can release the archive to free the compressed payload from the JavaScript heap, following prior work establishing that bundled, compressed delivery is critical to responsiveness on bandwidth-constrained networks~\cite{Evans2017}.

The per-frame sequence matches Unity's ordering: input polling, then zero or more fixed-timestep physics steps, then \texttt{update} across all enabled components followed by the active scenario, then \texttt{lateUpdate} in the same order, then transform synchronization, and finally rendering through the active camera before per-frame input state is reset.

\subsection{Deployment status}

The engine has been the 3D runtime of the KSU24 virtual learning environment at Kherson State University since April 2026 and is currently in \emph{closed testing}, with student access not yet opened while development and defect resolution continue. Ten educational scenarios have been built on it and packaged for the platform, spanning physics (mechanics, optics, electrical circuits, kinetic theory of gases, Kepler's laws), chemistry (atomic structure, molecular models, crystal lattices), biology (DNA structure), and astronomy (an interactive Solar System), alongside three internal scenarios for regression and performance testing. Nine of the ten educational scenarios are built procedurally from engine primitives and ship as archives of 10--20 KB, while the texture-based Solar System scenario is 17 MB.